\documentclass[11pt]{article}
\usepackage{amsmath, amssymb, geometry, amsthm, graphicx, xcolor}
\geometry{letterpaper, margin=1in}

\title{Homework 1.9: The Matrix of a Linear Transformation}
\author{Tashfeen Omran}
\date{September 2026}

\begin{document}

\maketitle
\tableofcontents
\newpage

\section{Part 1: Comprehensive Introduction, Context, and Prerequisites}

\subsection{Core Concepts}
The central idea of this chapter is that \textit{every} linear transformation $T: \mathbb{R}^n \to \mathbb{R}^m$ can be perfectly represented by a matrix multiplication. That is, for any such linear mapping, there exists a unique $m \times n$ matrix $A$ such that $T(\mathbf{x}) = A\mathbf{x}$ for all vectors $\mathbf{x}$ in $\mathbb{R}^n$. 

To find this matrix $A$ (called the \textbf{standard matrix} for the linear transformation), we only need to know what the transformation does to the standard basis vectors of $\mathbb{R}^n$. The standard basis vectors, denoted as $\mathbf{e}_1, \mathbf{e}_2, \dots, \mathbf{e}_n$, are the columns of the $n \times n$ identity matrix. If we know their images $T(\mathbf{e}_1), \dots, T(\mathbf{e}_n)$, we simply place these resulting vectors as the columns of our matrix $A$.

This section also introduces key terminology mapping back to functions:
\begin{itemize}
    \item \textbf{Onto (Surjective):} A mapping $T: \mathbb{R}^n \to \mathbb{R}^m$ is onto if every vector $\mathbf{b}$ in $\mathbb{R}^m$ is the image of \textit{at least one} vector $\mathbf{x}$ in $\mathbb{R}^n$.
    \item \textbf{One-to-One (Injective):} A mapping $T: \mathbb{R}^n \to \mathbb{R}^m$ is one-to-one if each vector $\mathbf{b}$ in $\mathbb{R}^m$ is the image of \textit{at most one} vector $\mathbf{x}$ in $\mathbb{R}^n$.
\end{itemize}

\subsection{Intuition and Derivation}
Why does the standard matrix theorem work? It stems directly from the definition of linearity. By definition, a linear transformation satisfies $T(c\mathbf{u} + d\mathbf{v}) = cT(\mathbf{u}) + dT(\mathbf{v})$. 
Any vector $\mathbf{x} \in \mathbb{R}^n$ can be written as a linear combination of the standard basis vectors:
\[ \mathbf{x} = x_1\mathbf{e}_1 + x_2\mathbf{e}_2 + \dots + x_n\mathbf{e}_n \]
Applying the linear transformation $T$ to both sides, we get:
\[ T(\mathbf{x}) = T(x_1\mathbf{e}_1 + \dots + x_n\mathbf{e}_n) = x_1T(\mathbf{e}_1) + \dots + x_nT(\mathbf{e}_n) \]
By the definition of matrix-vector multiplication, this linear combination of vectors is exactly the same as multiplying a matrix containing $T(\mathbf{e}_i)$ as its columns by the vector $\mathbf{x}$:
\[ T(\mathbf{x}) = \begin{bmatrix} | & | & & | \\ T(\mathbf{e}_1) & T(\mathbf{e}_2) & \dots & T(\mathbf{e}_n) \\ | & | & & | \end{bmatrix} \begin{bmatrix} x_1 \\ x_2 \\ \vdots \\ x_n \end{bmatrix} = A\mathbf{x} \]
This derivation proves that $T$ is completely determined by where it sends the basis vectors!

\subsection{Historical Context and Motivation}
The algebraic codification of transformations into matrices was formally developed in the mid-19th century by mathematicians like Arthur Cayley and James Joseph Sylvester. Mathematicians and physicists were dealing with complex systems of linear substitutions in projective geometry and celestial mechanics. Writing out system after system of equations with variables $x, y, z$ was tedious and error-prone. 

Cayley realized that the coefficients of these linear substitutions could be detached from the variables and manipulated as single entities (matrices). This meant that complex geometric concepts (like rotating an object in space, then reflecting it, then shearing it) could be reduced to simple, mechanized matrix multiplications, drastically accelerating the pace of modern physics and engineering.

\subsection{Key Formulas}
\begin{itemize}
    \item \textbf{The Standard Matrix:} $A = \begin{bmatrix} T(\mathbf{e}_1) & T(\mathbf{e}_2) & \dots & T(\mathbf{e}_n) \end{bmatrix}$
    \item \textbf{Rotation Matrix in $\mathbb{R}^2$:} To rotate points counterclockwise by an angle $\theta$:
    \[ A = \begin{bmatrix} \cos\theta & -\sin\theta \\ \sin\theta & \cos\theta \end{bmatrix} \]
    \item \textbf{Theorem 11 (One-to-One):} $T$ is one-to-one $\iff$ $T(\mathbf{x}) = \mathbf{0}$ has only the trivial solution ($\mathbf{x} = \mathbf{0}$) $\iff$ Columns of $A$ are linearly independent.
    \item \textbf{Theorem 12 (Onto):} $T$ maps $\mathbb{R}^n$ onto $\mathbb{R}^m$ $\iff$ Columns of $A$ span $\mathbb{R}^m$.
\end{itemize}

\subsection{Prerequisites}
To master this section, you must have a solid grasp of:
\begin{enumerate}
    \item \textbf{Matrix-Vector Multiplication:} Knowing how to compute $A\mathbf{x}$ using the linear combination of columns method.
    \item \textbf{Linear Independence and Span:} Understanding these definitions from previous chapters is non-negotiable for understanding one-to-one and onto mappings.
    \item \textbf{Pivot Positions:} Being able to row-reduce a matrix to identify pivot columns and rows.
    \item \textbf{Basic Trigonometry:} Knowing the exact values of sine and cosine for standard unit circle angles ($\pi/2, \pi/4, \pi/3$, etc.).
\end{enumerate}

\newpage
\section{Part 2: Detailed Homework Solutions}
\textit{Note: Per the requested strategy, a representative sample covering every distinct problem type from the exercises has been selected and solved step-by-step.}

\subsection*{Problem 1 (Constructing from Basis Vectors)}
\textbf{Question:} Assume $T$ is a linear transformation. Find the standard matrix of $T$. 
$T : \mathbb{R}^2 \to \mathbb{R}^4$, $T(\mathbf{e}_1) = (2, 1, 2, 1)$ and $T(\mathbf{e}_2) = (-5, 2, 0, 0)$, where $\mathbf{e}_1 = (1, 0)$ and $\mathbf{e}_2 = (0, 1)$.

\textbf{Solution:}
\begin{enumerate}
    \item \textbf{State the Theorem:} The standard matrix $A$ of a linear transformation $T: \mathbb{R}^n \to \mathbb{R}^m$ is formed by making the images of the standard basis vectors the columns of the matrix.
    \[ A = \begin{bmatrix} T(\mathbf{e}_1) & T(\mathbf{e}_2) \end{bmatrix} \]
    \item \textbf{Format the Vectors:} Write the given images as column vectors.
    \[ T(\mathbf{e}_1) = \begin{bmatrix} 2 \\ 1 \\ 2 \\ 1 \end{bmatrix}, \quad T(\mathbf{e}_2) = \begin{bmatrix} -5 \\ 2 \\ 0 \\ 0 \end{bmatrix} \]
    \item \textbf{Construct the Matrix:} Place these column vectors side-by-side to form the $4 \times 2$ matrix $A$.
    \[ A = \begin{bmatrix} 2 & -5 \\ 1 & 2 \\ 2 & 0 \\ 1 & 0 \end{bmatrix} \]
\end{enumerate}
\textbf{Final Answer:} \[ A = \begin{bmatrix} 2 & -5 \\ 1 & 2 \\ 2 & 0 \\ 1 & 0 \end{bmatrix} \]

\subsection*{Problem 3 (Pure Geometric Transformation)}
\textbf{Question:} Find the standard matrix of $T: \mathbb{R}^2 \to \mathbb{R}^2$ that rotates points (about the origin) through $3\pi/2$ radians (in the counterclockwise direction).

\textbf{Solution:}
\begin{enumerate}
    \item \textbf{Identify the Transformation:} This is a standard 2D rotation. The matrix for a counterclockwise rotation by $\theta$ is:
    \[ A = \begin{bmatrix} \cos\theta & -\sin\theta \\ \sin\theta & \cos\theta \end{bmatrix} \]
    \item \textbf{Evaluate the Trigonometric Functions:} Here, $\theta = 3\pi/2$. 
    \[ \cos(3\pi/2) = 0 \]
    \[ \sin(3\pi/2) = -1 \]
    \item \textbf{Substitute into the Matrix:}
    \[ A = \begin{bmatrix} 0 & -(-1) \\ -1 & 0 \end{bmatrix} = \begin{bmatrix} 0 & 1 \\ -1 & 0 \end{bmatrix} \]
\end{enumerate}
\textbf{Final Answer:} \[ A = \begin{bmatrix} 0 & 1 \\ -1 & 0 \end{bmatrix} \]

\subsection*{Problem 7 (Composite Geometric Transformation)}
\textbf{Question:} Find the standard matrix of $T: \mathbb{R}^2 \to \mathbb{R}^2$ that first rotates points through $-3\pi/4$ radians and then reflects points through the horizontal $x_1$-axis.

\textbf{Solution:}
\begin{enumerate}
    \item \textbf{Analyze the First Transformation (Rotation):} Let $R$ be the rotation matrix for $\theta = -3\pi/4$.
    \[ R = \begin{bmatrix} \cos(-3\pi/4) & -\sin(-3\pi/4) \\ \sin(-3\pi/4) & \cos(-3\pi/4) \end{bmatrix} \]
    Since $-3\pi/4$ is in the third quadrant, both sine and cosine are negative:
    \[ \cos(-3\pi/4) = -\frac{1}{\sqrt{2}}, \quad \sin(-3\pi/4) = -\frac{1}{\sqrt{2}} \]
    \[ R = \begin{bmatrix} -1/\sqrt{2} & -(-1/\sqrt{2}) \\ -1/\sqrt{2} & -1/\sqrt{2} \end{bmatrix} = \begin{bmatrix} -1/\sqrt{2} & 1/\sqrt{2} \\ -1/\sqrt{2} & -1/\sqrt{2} \end{bmatrix} \]
    \item \textbf{Analyze the Second Transformation (Reflection):} Let $F$ be the matrix for reflection across the $x_1$-axis. This transformation leaves $x_1$ unchanged but negates $x_2$. Thus, $T(\mathbf{e}_1) = \mathbf{e}_1$ and $T(\mathbf{e}_2) = -\mathbf{e}_2$.
    \[ F = \begin{bmatrix} 1 & 0 \\ 0 & -1 \end{bmatrix} \]
    \item \textbf{Compose the Transformations:} To "first rotate, then reflect," we evaluate $F(R\mathbf{x})$. By matrix multiplication associativity, this is $(FR)\mathbf{x}$. The standard matrix is $A = FR$. \textit{Note the right-to-left order!}
    \[ A = \begin{bmatrix} 1 & 0 \\ 0 & -1 \end{bmatrix} \begin{bmatrix} -1/\sqrt{2} & 1/\sqrt{2} \\ -1/\sqrt{2} & -1/\sqrt{2} \end{bmatrix} \]
    \item \textbf{Perform Matrix Multiplication:} Multiply rows of $F$ by columns of $R$.
    Row 1 of $F$ leaves Row 1 of $R$ unchanged. Row 2 of $F$ negates Row 2 of $R$.
    \[ A = \begin{bmatrix} -1/\sqrt{2} & 1/\sqrt{2} \\ 1/\sqrt{2} & 1/\sqrt{2} \end{bmatrix} \]
\end{enumerate}
\textbf{Final Answer:} \[ A = \begin{bmatrix} -1/\sqrt{2} & 1/\sqrt{2} \\ 1/\sqrt{2} & 1/\sqrt{2} \end{bmatrix} \]

\subsection*{Problem 17 (Matrix from Abstract Formula)}
\textbf{Question:} Show that $T$ is a linear transformation by finding a matrix that implements the mapping. 
$T(x_1, x_2, x_3, x_4) = (0, x_1+x_2, x_2+x_3, x_3+x_4)$.

\textbf{Solution:}
\begin{enumerate}
    \item \textbf{Identify Domain and Codomain:} The input has 4 variables, so domain is $\mathbb{R}^4$. The output is a 4-tuple, so codomain is $\mathbb{R}^4$. We are looking for a $4 \times 4$ matrix $A$.
    \item \textbf{Extract Coefficients:} Write the output as a column vector and express it as a matrix multiplying the input vector $\mathbf{x}$:
    \[ T(\mathbf{x}) = \begin{bmatrix} 0 \\ x_1 + x_2 \\ x_2 + x_3 \\ x_3 + x_4 \end{bmatrix} \]
    \item \textbf{Construct Matrix Columns (or read coefficients directly):}
    Equation 1 is $0 = 0x_1 + 0x_2 + 0x_3 + 0x_4$.
    Equation 2 is $1x_1 + 1x_2 + 0x_3 + 0x_4$.
    Equation 3 is $0x_1 + 1x_2 + 1x_3 + 0x_4$.
    Equation 4 is $0x_1 + 0x_2 + 1x_3 + 1x_4$.
    \[ T(\mathbf{x}) = \begin{bmatrix} 0 & 0 & 0 & 0 \\ 1 & 1 & 0 & 0 \\ 0 & 1 & 1 & 0 \\ 0 & 0 & 1 & 1 \end{bmatrix} \begin{bmatrix} x_1 \\ x_2 \\ x_3 \\ x_4 \end{bmatrix} \]
    \item \textbf{Conclusion:} Because we successfully wrote $T(\mathbf{x})$ in the form $A\mathbf{x}$, $T$ is a linear transformation.
\end{enumerate}
\textbf{Final Answer:} \[ A = \begin{bmatrix} 0 & 0 & 0 & 0 \\ 1 & 1 & 0 & 0 \\ 0 & 1 & 1 & 0 \\ 0 & 0 & 1 & 1 \end{bmatrix} \]

\subsection*{Problem 21 (Finding Pre-Image / Solving Equations)}
\textbf{Question:} Let $T: \mathbb{R}^2 \to \mathbb{R}^2$ be a linear transformation such that $T(x_1, x_2) = (x_1+x_2, 4x_1+5x_2)$. Find $\mathbf{x}$ such that $T(\mathbf{x}) = (3, 8)$.

\textbf{Solution:}
\begin{enumerate}
    \item \textbf{Find the Standard Matrix $A$:} Using the coefficients from the given formula:
    \[ A = \begin{bmatrix} 1 & 1 \\ 4 & 5 \end{bmatrix} \]
    \item \textbf{Set up the System:} We need to solve $A\mathbf{x} = \mathbf{b}$, where $\mathbf{b} = \begin{bmatrix} 3 \\ 8 \end{bmatrix}$.
    \[ \begin{bmatrix} 1 & 1 \\ 4 & 5 \end{bmatrix} \begin{bmatrix} x_1 \\ x_2 \end{bmatrix} = \begin{bmatrix} 3 \\ 8 \end{bmatrix} \]
    \item \textbf{Form the Augmented Matrix and Row Reduce:}
    \[ \left[ \begin{array}{cc|c} 1 & 1 & 3 \\ 4 & 5 & 8 \end{array} \right] \]
    Replace Row 2 with (Row 2 - $4 \times$ Row 1):
    \[ \left[ \begin{array}{cc|c} 1 & 1 & 3 \\ 0 & 1 & -4 \end{array} \right] \]
    Replace Row 1 with (Row 1 - Row 2):
    \[ \left[ \begin{array}{cc|c} 1 & 0 & 7 \\ 0 & 1 & -4 \end{array} \right] \]
    \item \textbf{Extract Solution:} The reduced row echelon form (RREF) tells us $x_1 = 7$ and $x_2 = -4$.
\end{enumerate}
\textbf{Final Answer:} $\mathbf{x} = \begin{bmatrix} 7 \\ -4 \end{bmatrix}$

\subsection*{Problem 23 (True/False Theory)}
\textbf{Question:} True or False: A linear transformation $T: \mathbb{R}^n \to \mathbb{R}^m$ is completely determined by its effect on the columns of the $n \times n$ identity matrix. Justify your answer.

\textbf{Solution:}
\textbf{True.} The columns of the $n \times n$ identity matrix are exactly the standard basis vectors $\mathbf{e}_1, \dots, \mathbf{e}_n$ for $\mathbb{R}^n$. Because any vector $\mathbf{x}$ in $\mathbb{R}^n$ can be written as a linear combination of these basis vectors, the linearity property ($T(c\mathbf{u}+d\mathbf{v}) = cT(\mathbf{u})+dT(\mathbf{v})$) guarantees that knowing $T(\mathbf{e}_i)$ is sufficient to compute $T(\mathbf{x})$ for any $\mathbf{x}$. This is the foundational proof for Theorem 10.

\subsection*{Problem 33 (Injectivity and Surjectivity Analysis)}
\textbf{Question:} Determine if the specified linear transformation in Exercise 17 is (a) one-to-one and (b) onto. Justify each answer.

\textbf{Solution:}
\begin{enumerate}
    \item \textbf{Recall the Matrix:} From Exercise 17, $T: \mathbb{R}^4 \to \mathbb{R}^4$ with standard matrix:
    \[ A = \begin{bmatrix} 0 & 0 & 0 & 0 \\ 1 & 1 & 0 & 0 \\ 0 & 1 & 1 & 0 \\ 0 & 0 & 1 & 1 \end{bmatrix} \]
    \item \textbf{Analyze Pivots:} Notice that the first row of $A$ consists entirely of zeros. This means that regardless of row reduction, there can be at most 3 pivot positions in this $4 \times 4$ matrix. 
    \item \textbf{Test for Onto:} By Theorem 12, $T$ maps $\mathbb{R}^n$ onto $\mathbb{R}^m$ if and only if the columns of $A$ span $\mathbb{R}^m$, which requires a pivot in every row. Since row 1 will never have a pivot, $A$ does not have a pivot in every row. Thus, $T$ is \textbf{not onto}.
    \item \textbf{Test for One-to-One:} By Theorem 12, $T$ is one-to-one if and only if the columns of $A$ are linearly independent, which requires a pivot in every column. Since there are at most 3 pivots and 4 columns, there will be at least one free variable. Thus, $T$ is \textbf{not one-to-one}.
\end{enumerate}
\textbf{Final Answer:} (a) Not one-to-one. (b) Not onto.

\subsection*{Problem 39 (Conceptual Linking)}
\textbf{Question:} Let $T: \mathbb{R}^n \to \mathbb{R}^m$ be a linear transformation, with $A$ its standard matrix. Complete the following statement to make it true: "$T$ is one-to-one if and only if $A$ has \_\_\_\_ pivot columns." Explain why the statement is true.

\textbf{Solution:}
\textbf{Final Answer:} $T$ is one-to-one if and only if $A$ has \textbf{$n$} pivot columns. 
\textit{Explanation:} $T$ is one-to-one if and only if the equation $T(\mathbf{x}) = \mathbf{0}$ (which is $A\mathbf{x} = \mathbf{0}$) has only the trivial solution. This occurs if and only if there are no free variables in the system $A\mathbf{x} = \mathbf{0}$. For an $m \times n$ matrix $A$ to have no free variables, every one of its $n$ columns must be a pivot column. Therefore, $A$ must have exactly $n$ pivot columns.

\newpage
\section{Part 3: In-Depth Analysis of Problems and Techniques}

\subsection{A) Problem Types and General Approach}
We can categorize the homework into four main operational types:

\begin{enumerate}
    \item \textbf{Matrix Construction from Basis (Type 1):} e.g., Problem 1.
    \textit{Approach:} The easiest problem type. You are handed the columns of the matrix. Simply place $T(\mathbf{e}_1), T(\mathbf{e}_2), \dots$ into the columns of a new matrix in ascending order. No calculation is needed, just direct transcription.
    
    \item \textbf{Geometric Transformations (Type 2):} e.g., Problems 3, 7.
    \textit{Approach:} These require translating a visual or verbal geometric action into a matrix. If a standard formula is known (like rotation), plug the angle $\theta$ into the rotation matrix. If multiple transformations occur, create separate matrices and multiply them. \textbf{Crucial Strategy:} Matrix multiplication must be written Right-to-Left relative to the sequence of events (i.e., transformation 1 is the rightmost matrix).
    
    \item \textbf{Formula to Matrix Translation (Type 3):} e.g., Problem 17, 21.
    \textit{Approach:} You are given an abstract algebraic formula. Read off the coefficients of $x_1, x_2, \dots$ to construct the matrix row by row. This is effectively reversing matrix-vector multiplication.
    
    \item \textbf{Onto and One-to-One Analysis (Type 4):} e.g., Problems 33, 39.
    \textit{Approach:} Always boil these concepts down to \textbf{pivot positions}. Form the matrix $A$, row reduce it, and count the pivots. 
    \begin{itemize}
        \item Does every column have a pivot? Yes $\implies$ One-to-One.
        \item Does every row have a pivot? Yes $\implies$ Onto.
    \end{itemize}
\end{enumerate}

\subsection{B) Key Algebraic and Linear Algebra Manipulations}
\begin{itemize}
    \item \textbf{Evaluating Trigonometric Boundary Values:} In Problem 3, simplifying $\cos(3\pi/2)$ and $\sin(3\pi/2)$ was required. Recognizing quadrant signs correctly is the most common pitfall here.
    \item \textbf{Right-to-Left Matrix Multiplication for Composition:} In Problem 7, applying a rotation $R$ \textit{then} reflection $F$ required computing the product $FR$, not $RF$. Because $T(\mathbf{x}) = F(R(\mathbf{x})) = (FR)\mathbf{x}$. Left-multiplication preserves function composition order.
    \item \textbf{Solving Linear Systems via Augmented Matrices:} In Problem 21, the task was to find the input vector $\mathbf{x}$ that produced a known output. This transforms abstract function evaluation into a concrete linear system $A\mathbf{x} = \mathbf{b}$, requiring classical row reduction to Reduced Row Echelon Form (RREF).
    \item \textbf{Zero-Row Pivot Analysis:} In Problem 33, observing that the unreduced matrix had a row of pure zeros instantly bypassed the need to do full row reduction. A row of zeros means a maximum of $m-1$ pivots, immediately proving the $m \times n$ matrix cannot have a pivot in every row.
\end{itemize}

\newpage
\section{Part 4: "Cheatsheet" and Tips for Success}

\subsection{Important Formulas}
\begin{itemize}
    \item \textbf{Standard Matrix Definition:} $A = \begin{bmatrix} T(\mathbf{e}_1) & T(\mathbf{e}_2) & \dots & T(\mathbf{e}_n) \end{bmatrix}$
    \item \textbf{2D Rotation Matrix (Counterclockwise by $\theta$):} 
    \[ R_\theta = \begin{bmatrix} \cos\theta & -\sin\theta \\ \sin\theta & \cos\theta \end{bmatrix} \]
    \item \textbf{2D Reflection Matrices:}
    \begin{itemize}
        \item Across $x_1$-axis ($y=0$): $\begin{bmatrix} 1 & 0 \\ 0 & -1 \end{bmatrix}$
        \item Across $x_2$-axis ($x=0$): $\begin{bmatrix} -1 & 0 \\ 0 & 1 \end{bmatrix}$
        \item Across $x_2=x_1$ ($y=x$): $\begin{bmatrix} 0 & 1 \\ 1 & 0 \end{bmatrix}$
    \end{itemize}
\end{itemize}

\subsection{Tricks and "Cheats"}
\begin{itemize}
    \item \textbf{The $T(\mathbf{e}_i)$ "Cheat":} If you forget the standard formula for a reflection or shear, don't panic. Just draw a quick Cartesian plane. Where does the point $(1,0)$ go? Where does $(0,1)$ go? Those two new coordinates form the first and second columns of your matrix, respectively. You can derive ANY 2D geometric matrix in 15 seconds this way.
    \item \textbf{The Dimension Check:} The dimensions of the transformation give you the matrix size immediately. $T: \mathbb{R}^n \to \mathbb{R}^m$ yields an $m \times n$ matrix ($m$ rows, $n$ columns). Memorize this phrase: \textit{"Outputs by Inputs"}.
\end{itemize}

\subsection{Common Pitfalls and Mistakes}
\begin{itemize}
    \item \textbf{The Composition Trap:} If a problem states "Transform A, THEN Transform B", students almost universally write $A \times B$. It must be $B \times A$. Matrix multiplication is evaluated from the vector outward (right to left). 
    \item \textbf{Clockwise vs. Counterclockwise:} The standard rotation matrix formula assumes a positive, counterclockwise angle. If a problem states "rotate clockwise by $90^\circ$", you MUST plug in $\theta = -90^\circ$ (or $-\pi/2$) into the formula.
    \item \textbf{Confusing One-to-One and Onto Pivot Rules:}
    \begin{itemize}
        \item Pivot in every \textbf{COLUMN} $\to$ \textbf{C}ounting independent inputs $\to$ One-to-One.
        \item Pivot in every \textbf{ROW} $\to$ \textbf{R}eaching all outputs $\to$ Onto.
    \end{itemize}
\end{itemize}

\newpage
\section{Part 5: Conceptual Synthesis and The "Big Picture"}

\subsection{A) Thematic Connections}
The dominant theme in this chapter is \textbf{Representation and Isomorphism}. Mathematics frequently involves shifting a problem from a difficult domain into a simpler one. Here, we shift the abstract domain of "functions and mappings" into the concrete, computational domain of "tables of numbers (matrices)". 

This perfectly mirrors the shift in calculus where we use Taylor Series to turn complex transcendental functions (like $e^x$ or $\sin(x)$) into basic polynomials. Matrices do the same thing for linear functions; they provide a computational chassis that allows a computer to perfectly model and execute functional mappings.

\subsection{B) Forward and Backward Links}
\textbf{Backward Links:} This topic is the culmination of Chapter 1. It bridges the gap between solving systems of linear equations ($A\mathbf{x} = \mathbf{b}$) and understanding vectors algebraically. The pivot theorems for 1-to-1 and onto map directly backward to the rules for determining if a linear system has a unique solution, infinite solutions, or no solution.

\textbf{Forward Links (The Bedrock for Future Math):}
\begin{itemize}
    \item \textbf{Computer Graphics:} The geometric transformation matrices learned here (rotations, shears) are the literal engines behind rendering 3D graphics in computer science.
    \item \textbf{Multivariable Calculus (The Jacobian):} Later in calculus, you will approximate \textit{non-linear} multivariable functions using linear transformations. The matrix that represents this "best linear approximation" is called the Jacobian Matrix, and understanding how it maps infinitesimally small vectors relies completely on the foundation built in Section 1.9.
\end{itemize}

\newpage
\section{Part 6: Real-World Application and Modeling (Finance Focus)}

\subsection{A) Concrete Scenarios}
\begin{enumerate}
    \item \textbf{Portfolio Risk Factor Mapping:} In quantitative finance, the returns of thousands of stocks are often modeled as linear combinations of a smaller set of underlying "macro-economic risk factors" (like interest rates, inflation, market momentum). A risk manager maps a vector of daily changes in these abstract factors ($\mathbb{R}^n$) into a vector of daily price changes for their specific portfolio of assets ($\mathbb{R}^m$). This mapping is a linear transformation encoded by a factor loading matrix.
    \item \textbf{Transition Matrices in Credit Ratings:} The probability that a bond gets upgraded (e.g., from AA to AAA) or downgraded into default over a year is modeled using linear mappings. A vector representing the current distribution of a bank's debt is multiplied by a transition matrix (representing the mapping transformation) to predict the future risk profile of the portfolio.
    \item \textbf{Hedging and Replicating Portfolios:} In options pricing (derivatives), traders construct "replicating portfolios" using stocks and risk-free bonds to perfectly mimic the payoff of a complex derivative. Determining the exact quantity of stock and bonds to buy requires solving a linear transformation mapping the asset quantities to the future payoff states.
\end{enumerate}

\subsection{B) Model Problem Setup: Portfolio Risk Factor Mapping}
\textbf{Scenario:} A quantitative hedge fund manages a portfolio of $m = 500$ different stocks. To manage risk, they map the returns of these stocks to a set of $n = 3$ core macroeconomic factors: Interest Rate Risk, Market Risk, and Sector Risk.

\textbf{Variables:}
\begin{itemize}
    \item Let $\mathbf{f} \in \mathbb{R}^3$ be the vector representing the daily percentage change in the 3 risk factors.
    \item Let $\mathbf{r} \in \mathbb{R}^{500}$ be the vector representing the daily percentage return of the 500 stocks.
    \item Let $A$ be the $500 \times 3$ \textit{Factor Loading Matrix}. Each column $j$ of $A$ represents the standard basis mapping $T(\mathbf{e}_j)$: how the entire portfolio of 500 stocks reacts if factor $j$ moves by 1 unit while the others remain zero.
\end{itemize}

\textbf{Mathematical Model:}
The linear mapping from the factor-space to the stock-return-space is defined by the transformation $T(\mathbf{f}) = A\mathbf{f}$.
If the portfolio manager wants to construct a portfolio that is perfectly "delta-neutral" (zero change in returns) to a specific predicted market shock $\mathbf{f}_{shock}$, they need to solve the linear equation:
\[ A\mathbf{f}_{shock} + \mathbf{w}_{hedge} = \mathbf{0} \]
Where they must calculate the necessary hedge vector $\mathbf{w}_{hedge}$ in $\mathbb{R}^{500}$ to offset the linear transformation of the factor shock. The pivot structure of $A$ dictates if the fund has enough independent assets to perfectly span the risk and achieve this hedge (a direct application of the "Onto" property).

\newpage
\section{Part 7: Common Variations and Untested Concepts}

While the homework covered basis vectors, formulas, and geometry, linear algebra exams frequently test the boundaries of these definitions. Here is an important variation not explicitly in your assignment.

\subsection{Variation: Constructing a Matrix from NON-Basis Vectors}
\textbf{Concept:} What if you are given the images of random vectors, not the standard basis vectors $\mathbf{e}_1$ and $\mathbf{e}_2$? You cannot simply place the given images into the columns of a matrix. You must use the linearity property to solve for $T(\mathbf{e}_1)$ and $T(\mathbf{e}_2)$ first.

\textbf{Example Problem:}
Suppose $T: \mathbb{R}^2 \to \mathbb{R}^2$ is a linear transformation. You are given that $T\left( \begin{bmatrix} 1 \\ 1 \end{bmatrix} \right) = \begin{bmatrix} 3 \\ 5 \end{bmatrix}$ and $T\left( \begin{bmatrix} 1 \\ -1 \end{bmatrix} \right) = \begin{bmatrix} 1 \\ 1 \end{bmatrix}$. Find the standard matrix $A$ for $T$.

\textbf{Solution:}
We need $T(\mathbf{e}_1)$ and $T(\mathbf{e}_2)$. 
Notice that we can express $\mathbf{e}_1$ as a linear combination of the two given input vectors:
\[ \mathbf{e}_1 = \begin{bmatrix} 1 \\ 0 \end{bmatrix} = \frac{1}{2} \begin{bmatrix} 1 \\ 1 \end{bmatrix} + \frac{1}{2} \begin{bmatrix} 1 \\ -1 \end{bmatrix} \]
Using the linearity property $T(c\mathbf{u} + d\mathbf{v}) = cT(\mathbf{u}) + dT(\mathbf{v})$:
\begin{align*}
    T(\mathbf{e}_1) &= \frac{1}{2} T\left( \begin{bmatrix} 1 \\ 1 \end{bmatrix} \right) + \frac{1}{2} T\left( \begin{bmatrix} 1 \\ -1 \end{bmatrix} \right) \\
    &= \frac{1}{2} \begin{bmatrix} 3 \\ 5 \end{bmatrix} + \frac{1}{2} \begin{bmatrix} 1 \\ 1 \end{bmatrix} = \begin{bmatrix} 1.5 \\ 2.5 \end{bmatrix} + \begin{bmatrix} 0.5 \\ 0.5 \end{bmatrix} = \begin{bmatrix} 2 \\ 3 \end{bmatrix}
\end{align*}
Similarly, for $\mathbf{e}_2$:
\[ \mathbf{e}_2 = \begin{bmatrix} 0 \\ 1 \end{bmatrix} = \frac{1}{2} \begin{bmatrix} 1 \\ 1 \end{bmatrix} - \frac{1}{2} \begin{bmatrix} 1 \\ -1 \end{bmatrix} \]
\begin{align*}
    T(\mathbf{e}_2) &= \frac{1}{2} \begin{bmatrix} 3 \\ 5 \end{bmatrix} - \frac{1}{2} \begin{bmatrix} 1 \\ 1 \end{bmatrix} = \begin{bmatrix} 1.5 \\ 2.5 \end{bmatrix} - \begin{bmatrix} 0.5 \\ 0.5 \end{bmatrix} = \begin{bmatrix} 1 \\ 2 \end{bmatrix}
\end{align*}
Now that we have the images of the standard basis vectors, we construct the matrix:
\[ A = \begin{bmatrix} T(\mathbf{e}_1) & T(\mathbf{e}_2) \end{bmatrix} = \begin{bmatrix} 2 & 1 \\ 3 & 2 \end{bmatrix} \]

\newpage
\section{Part 8: Advanced Diagnostic Testing: "Find the Flaw"}
\textit{Instructions: Each of the following solutions contains a single, critical conceptual or algebraic error. Identify the error, explain why it is wrong, and provide the correct final answer.}

\subsection*{Flawed Problem 1}
\textbf{Question:} Let $T: \mathbb{R}^2 \to \mathbb{R}^2$ be a linear transformation where $T(\mathbf{e}_1) = (4, -1)$ and $T(\mathbf{e}_2) = (2, 7)$. Find the standard matrix $A$.
\textbf{Flawed Solution:} 
By the standard matrix theorem, we group the outputs into the matrix $A$.
\[ A = \begin{bmatrix} 4 & -1 \\ 2 & 7 \end{bmatrix} \]

\vspace{2cm}
\subsection*{Flawed Problem 2}
\textbf{Question:} Let $T$ be a transformation whose standard matrix is $A = \begin{bmatrix} 1 & 3 & 5 \\ 0 & 2 & 4 \end{bmatrix}$. Determine if $T$ is a one-to-one mapping.
\textbf{Flawed Solution:} 
We look at the pivots of the matrix. 
\[ A = \begin{bmatrix} 1 & 3 & 5 \\ 0 & 2 & 4 \end{bmatrix} \]
There is a pivot in the first row (1) and a pivot in the second row (2). Because the matrix has a pivot in every row, the mapping $T$ is one-to-one.

\vspace{2cm}
\subsection*{Flawed Problem 3}
\textbf{Question:} Find the standard matrix for the transformation $T: \mathbb{R}^2 \to \mathbb{R}^2$ that reflects points through the line $x_2 = -x_1$.
\textbf{Flawed Solution:}
The line $x_2 = -x_1$ is the same as the line $y = x$ but with a negative sign. The matrix for a reflection across $y = x$ flips the coordinates, so $x_1$ becomes $x_2$ and $x_2$ becomes $x_1$.
\[ A = \begin{bmatrix} 0 & 1 \\ 1 & 0 \end{bmatrix} \]

\vspace{2cm}
\subsection*{Flawed Problem 4}
\textbf{Question:} Find the standard matrix for $T: \mathbb{R}^2 \to \mathbb{R}^2$ that first performs a rotation of $\pi/2$ radians counterclockwise, and then reflects points across the $x_1$-axis.
\textbf{Flawed Solution:}
Let $R$ be the rotation matrix: $R = \begin{bmatrix} 0 & -1 \\ 1 & 0 \end{bmatrix}$.
Let $F$ be the reflection matrix across $x_1$: $F = \begin{bmatrix} 1 & 0 \\ 0 & -1 \end{bmatrix}$.
We multiply them in order to find the combined transformation matrix $A$:
\[ A = R \times F = \begin{bmatrix} 0 & -1 \\ 1 & 0 \end{bmatrix} \begin{bmatrix} 1 & 0 \\ 0 & -1 \end{bmatrix} = \begin{bmatrix} 0 & 1 \\ 1 & 0 \end{bmatrix} \]

\vspace{2cm}
\subsection*{Flawed Problem 5}
\textbf{Question:} Find the standard matrix for the transformation $T: \mathbb{R}^2 \to \mathbb{R}^2$ defined by $T(x_1, x_2) = (x_1 + 3, x_2)$.
\textbf{Flawed Solution:}
We evaluate the standard basis vectors:
$T(\mathbf{e}_1) = T(1, 0) = (1 + 3, 0) = (4, 0)$.
$T(\mathbf{e}_2) = T(0, 1) = (0 + 3, 1) = (3, 1)$.
We construct the matrix using these columns:
\[ A = \begin{bmatrix} 4 & 3 \\ 0 & 1 \end{bmatrix} \]


\newpage
\subsection*{Answer Key to "Find the Flaw" Diagnostic}

\textbf{Flaw 1 Correction:}
\begin{itemize}
    \item \textbf{The Error:} The student placed the vectors as \textit{rows} instead of \textit{columns}. 
    \item \textbf{Correct Step:} The outputs $T(\mathbf{e}_1)$ and $T(\mathbf{e}_2)$ must form the columns of the standard matrix.
    \item \textbf{Correct Final Answer:} $A = \begin{bmatrix} 4 & 2 \\ -1 & 7 \end{bmatrix}$
\end{itemize}

\textbf{Flaw 2 Correction:}
\begin{itemize}
    \item \textbf{The Error:} The student confused the condition for "onto" with the condition for "one-to-one". A pivot in every row proves it is \textit{onto}, not one-to-one.
    \item \textbf{Correct Step:} For a mapping to be one-to-one, the matrix must have a pivot in every \textit{column}. This $2 \times 3$ matrix has 3 columns but only 2 pivots, leaving one free variable.
    \item \textbf{Correct Final Answer:} $T$ is NOT one-to-one.
\end{itemize}

\textbf{Flaw 3 Correction:}
\begin{itemize}
    \item \textbf{The Error:} The student used the matrix for reflection across $y=x$, but failed to apply the negative transformation for $y=-x$. 
    \item \textbf{Correct Step:} A reflection across $x_2 = -x_1$ maps $(1,0) \to (0,-1)$ and $(0,1) \to (-1,0)$. 
    \item \textbf{Correct Final Answer:} $A = \begin{bmatrix} 0 & -1 \\ -1 & 0 \end{bmatrix}$
\end{itemize}

\textbf{Flaw 4 Correction:}
\begin{itemize}
    \item \textbf{The Error:} The student multiplied the matrices left-to-right based on the english wording of the sentence. Matrix composition must be evaluated right-to-left. 
    \item \textbf{Correct Step:} Evaluate $F(R(\mathbf{x}))$, which means multiplying $F \times R$.
    \item \textbf{Correct Final Answer:} 
    \[ A = F \times R = \begin{bmatrix} 1 & 0 \\ 0 & -1 \end{bmatrix} \begin{bmatrix} 0 & -1 \\ 1 & 0 \end{bmatrix} = \begin{bmatrix} 0 & -1 \\ -1 & 0 \end{bmatrix} \]
\end{itemize}

\textbf{Flaw 5 Correction:}
\begin{itemize}
    \item \textbf{The Error:} The student assumed the transformation was linear and blindly applied the theorem. However, a translation (adding a constant $+3$) is NOT a linear transformation (e.g., $T(0,0) = (3,0) \neq (0,0)$).
    \item \textbf{Correct Step:} A linear transformation must satisfy $T(\mathbf{0}) = \mathbf{0}$. Since it does not, no standard matrix representation exists in $\mathbb{R}^2$.
    \item \textbf{Correct Final Answer:} $T$ is not a linear transformation, so no such $2 \times 2$ standard matrix exists.
\end{itemize}

\end{document}